\documentclass[10pt]{article}

% Public supplement to: Paired-Acquisition Neural Factorization: An End-to-End Computational Pathology Pipeline

\usepackage[letterpaper,top=0.48in,bottom=0.56in,left=0.62in,right=0.62in]{geometry}
\usepackage{amsmath,amssymb,amsfonts}
\usepackage{booktabs}
\usepackage{graphicx}
\usepackage{xcolor}
\usepackage{microtype}
\usepackage{placeins}
\usepackage{caption}
\usepackage{enumitem}
\usepackage{times}
\usepackage{titlesec}
\usepackage{array}
\usepackage{tabularx}
\usepackage{xurl}
\usepackage{hyperref}
\usepackage[numbers,sort&compress]{natbib}

\hypersetup{colorlinks=true,linkcolor=blue!55!black,citecolor=blue!55!black,urlcolor=blue!55!black}
\setlength{\parindent}{0pt}
\setlength{\parskip}{0.05em}
\setlength{\textfloatsep}{0.38em plus 0.15em minus 0.15em}
\setlength{\floatsep}{0.34em plus 0.12em minus 0.12em}
\setlength{\intextsep}{0.34em plus 0.12em minus 0.12em}
\captionsetup{font=small,labelfont=bf,skip=2pt}
\titlespacing*{\section}{0pt}{0.62em}{0.25em}
\newcolumntype{Y}{>{\raggedright\arraybackslash}X}
\newcommand{\codepath}[1]{{\ttfamily\footnotesize\path{#1}}}
\emergencystretch=3em

\begin{document}
\begin{center}
{\LARGE\bfseries Supplement --- Paired-Acquisition Neural Factorization: An End-to-End Computational Pathology Pipeline\par}
\vspace{0.04in}
{\small Public preprint supplement --- 2026-09-04. Research use only; not clinical software.}
\vspace{0.04in}
\rule{\textwidth}{1.0pt}
\end{center}

\section{Implementation-path inventory}
Table~\ref{tab:impl} lists the primary implementation paths per research line,
bound in the evidence manifest
(\codepath{evidence/manuscript_evidence_manifest.json}). Every architectural
claim in the main text binds to these source paths and their tests.

\begin{table}[!htbp]
\centering
\scriptsize
\caption{Implementation-path inventory.}
\label{tab:impl}
\begin{tabularx}{\textwidth}{@{}p{0.15\textwidth}Y@{}}
\toprule
Line & Primary implementation paths \\
\midrule
PCam & \codepath{experiments/train_pcam.py}; \codepath{evaluate_pcam.py}; \codepath{src/data/datasets/pcam_dataset.py} \\
PA-NF & \codepath{experiments/paired_acquisition/} runners; \codepath{src/paired_acquisition_factorial.py} \\
SCORPION & \codepath{experiments/paired_acquisition/run_*scorpion*}; \codepath{src/models/scorpion_pathoalign.py} \\
Canine & \codepath{run_biological_label_preservation_fixed_estimand.py} \\
TransnnMIL & \codepath{src/models/transnnmil/}; \codepath{src/models/mil/} \\
PathologyFL & \codepath{src/features/federated/pathology_fl/} \\
FAIR-WEIGHTS-H & \codepath{.../weighting/fair_weights_h.py}; \codepath{.../pathoalign_bridge.py} \\
CAMELYON17 & \codepath{experiments/camelyon17_*}; \codepath{scripts/camelyon17/} \\
Reproducibility & \codepath{scripts/provenance/}; \codepath{src/paired_acquisition_provenance.py} \\
\bottomrule
\end{tabularx}
\end{table}
\FloatBarrier

\section{Test coverage}
Table~\ref{tab:tests} maps each line to its test suites. Tests are the binding
target for architectural claims.

\begin{table}[!htbp]
\centering
\scriptsize
\caption{Test coverage per research line.}
\label{tab:tests}
\begin{tabularx}{\textwidth}{@{}p{0.22\textwidth}Y@{}}
\toprule
Line & Test suites \\
\midrule
PCam & \codepath{tests/test_pcam_*.py}; \codepath{test_pcam_evaluation_ci.py} \\
Paired/SCORPION/canine & \codepath{tests/test_fixed_estimand_real_feature_space_adjudication*.py}; \codepath{tests/test_real_paired_scanner_*.py} \\
Synthetic & \codepath{tests/test_biological_bottleneck_capacity_allocation_factorial.py} \\
TransnnMIL & \codepath{tests/models/test_transnnmil_*.py}; \codepath{test_topology_branch.py}; \codepath{test_hierarchical_pooling.py} \\
PathologyFL & \codepath{tests/federated/test_fl_integration.py}; \codepath{test_pathology_fl_*.py}; \codepath{test_secure_aggregation*.py} \\
FAIR-WEIGHTS-H & \codepath{tests/federated/test_fair_weights_h.py}; \codepath{test_weighting_*.py} \\
CAMELYON17 & \codepath{tests/test_camelyon_*.py}; \codepath{tests/camelyon17/} \\
Reproducibility & \codepath{tests/test_corrected_paired_acquisition_evidence.py}; \codepath{test_real_bottleneck_representation_recovery.py} \\
\bottomrule
\end{tabularx}
\end{table}
\FloatBarrier

\section{Result artifact hashes}
The authoritative artifact hashes are recorded in
\codepath{evidence/manuscript_evidence_manifest.json}. Selected current hashes:

\begin{itemize}
\item corrected paired-acquisition release manifest: \codepath{evidence/paired_acquisition/corrected-20260726/release_manifest.json};
\item SCORPION capacity-matched release manifest: \codepath{evidence/paired_acquisition/scorpion-capacity-matched-20260726/release_manifest.json};
\item dimensionality/cross-covariance release manifest: \codepath{evidence/paired_acquisition/dimensionality-xcov-factorial-20260726/release_manifest.json}.
\end{itemize}

\section{Original PA-NF objective and checkpoint provenance}
The registered 175-fit campaign was executed at training commit
\codepath{0adea50f1ef22865969109f1834a3c175e3f8b43}. Its original model
definition \codepath{src/models/scorpion_pathoalign.py} implements the objective
shown in Equation~\ref{eq:panf} of the main manuscript: a multi-positive
supervised-contrastive loss using matched region identities, \emph{same-view}
feature reconstruction by a decoder consuming concatenated B/A branches,
variance and covariance penalties, cross-branch covariance, normalized
B--scanner-label dependence, a gradient-reversed B scanner classifier, and
an ordinary A scanner classifier. It does not implement a crossed-target
decoder loss or a prototype loss in the released 25-checkpoint
\codepath{pathoalign_dep20} family. These details were checked against the exact
model source blob at the training commit and the 25-checkpoint model-release
audit. The release audit binds fold-specific feature standardization, 75
training epochs, checkpoint dimensions $768\to B256/A64$, and all 25
registered fold/seed checkpoints. Describing the losses as constraints on
measured information accessibility does not assume latent-factor uniqueness.

\section{PA-NF controlled-result boundary}
The primary registered SCORPION comparison is the 175-fit capacity-matched campaign: seven variants, five folds, and five seeds. PA-NF reduces tissue-branch scanner balanced accuracy by 0.3108 relative to the equal-capacity two-branch neural control while satisfying the registered same-region retrieval noninferiority criterion and retaining strong acquisition-branch scanner information.

The corrected canine SCC comparison answers a broader and harder question. Under the fixed five-category estimand, strong centroid/QR and paired-linear corrections remain competitive and no universal neural feature-space increment is established. The two results use different comparators and endpoints and should be interpreted together rather than collapsed into one ranking.

\section{Completed repaired TransnnMIL matched comparison}
A separate, frozen 35-fit PANDA comparison used seven architectures and five matched seeds with slide-disjoint training, model-selection, and confirmation subsets. The repaired canonical TransnnMIL mean confirmation QWK was 0.792195, versus 0.805826 for TransMIL and 0.779618 for nnMIL. The five matched repaired-minus-TransMIL QWK differences were negative (mean $-0.013631$); the registered primary architecture-superiority contrast failed. The narrower repaired-minus-nnMIL mean was $+0.012577$ (five of five seeds favorable). This study must not be conflated with the 18-run pre-remediation optimization grid. The original 35-cell checkpoint provenance was recorded on a separate September 2026 research branch; a later external-spec amendment integrity failure does not invalidate the recorded PANDA confirmation comparison but does block any stronger claim about the pending external SICAP run.

\section{Exploratory ScanComp three-view external retrieval}
Two independent physical slides, \texttt{kwo930} and \texttt{jye249}, each supplied 32 frozen held-out tissue regions observed under HMS360, HSTP1000 and LCAP450. Source-only frozen registration and original held-out geometry gates admitted both required scanner pairs, but PHUFS failed full four-scanner qualification on both slides. The subsequent same-region retrieval was therefore an exploratory post-QC-selected three-view study, not four-view prospective confirmation.

The original 25 SCORPION-trained PA-NF full B256/A64 checkpoints and 50 fold/seed-matched historical objective controls were evaluated with pinned DINOv2 X768 preprocessing and standardizers. Frozen full B256 improved mean strict top-1 from 92.96875\% to 97.50000\% on \texttt{kwo930} and from 81.25000\% to 88.21875\% on \texttt{jye249}. The full-minus-no-adversary mean cosine margins were 0.003301974 and 0.004705958; full-minus-two-branch-no-scanner-objectives margins were 0.012552280 and 0.010453726. These control comparisons use identical checkpoint families and are not independent specimens. The two-slide equal-weight strict-top-1 improvement is 5.75000 percentage points; no physical-slide-level population confidence interval follows from two distinct slides.

External scanner-variance decompositions did not establish suppression versus raw X768, and no independent human diagnostic ground truth is admitted. Additional source-selected registration-ineligible attempts must remain visible in source-eligibility accounting. Local source TIFF custody, original registration receipts, preselected slide ranks, per-query scores and frozen checkpoint hashes require separate preservation and release checking before any public reproducibility bundle claims that they are downloadable. No original negative or failed registration may be promoted to a positive outcome.

\section{Pre-remediation whole-slide optimization records}
The original pre-remediation PANDA execution path was evaluated before correction of its branch-fusion and topology mechanisms. These results remain part of the research record but are not estimates of performance for the repaired architecture.

This section preserves the earlier execution-path evidence. The corrected canonical implementation was subsequently evaluated in a completed, separately frozen seven-architecture by five-seed PANDA comparison (35 fits). The numerical 18-run stability grid below remains historical pre-remediation evidence and must not be confused with those repaired-model outcomes.

\subsection{Full-PANDA feature pipeline}
PANDA provides the primary slide-level substrate. The feature manifest was verified before training: 10,614 candidate HDF5 files were attempted, 10,611 were readable, and three unreadable files were removed. The resulting split contains 8,488 training slides and 2,123 validation slides, with 768-dimensional Phikon patch embeddings.

Before the later TransnnMIL implementation remediation, mean-pooled Phikon followed by an MLP reached validation QWK 0.7274, while gated AttentionMIL reached 0.8100. Three historical TransnnMIL executions reached 0.8155, 0.8225, and 0.8086. These values remain records of the historical execution path; because that path was later found to contain fusion and topology implementation defects, they do not validate genuine repaired TransMIL--nnMIL fusion or a trained repaired topology contribution.

\begin{table}[t]
\centering
\caption{Historical pre-remediation PANDA whole-slide execution results on Phikon feature bags. QWK is quadratic weighted kappa. The TransnnMIL values do not constitute validation of the repaired canonical architecture.}
\label{tab:panda_mil}
\small
\begin{tabular}{lr}
\toprule
Model / historical run & Validation QWK \\
\midrule
Mean-pooled Phikon + MLP & 0.7274 \\
Gated AttentionMIL & 0.8100 \\
Historical TransnnMIL seed 42 & 0.8155 \\
Historical TransnnMIL seed 123 & 0.8225 \\
Historical TransnnMIL seed 2025 & 0.8086 \\
\bottomrule
\end{tabular}
\end{table}

\subsection{Historical optimizer-stability record}
The pre-remediation TransnnMIL execution path exposed substantial learning-rate sensitivity. A stabilization recipe combining AdamW, warmup-cosine scheduling, two warmup epochs, gradient clipping at norm 1.0, and early stopping was evaluated across six learning rates and seeds 42, 123, and 2025, producing 18 historical full-PANDA runs.

\begin{table*}[t]
\centering
\caption{Historical pre-remediation TransnnMIL optimization-stability grid on the full readable PANDA feature manifest. These results characterize the earlier execution path and are not evidence for the repaired canonical implementation.}
\label{tab:transnnmil_stability}
\small
\begin{tabular}{lrrrr}
\toprule
Learning rate & Mean best val QWK & Std. & Min. & Max. \\
\midrule
$1\times10^{-4}$ & \textbf{0.8257} & 0.0169 & 0.8087 & 0.8425 \\
$2\times10^{-4}$ & 0.8245 & 0.0192 & 0.8077 & \textbf{0.8455} \\
$3\times10^{-4}$ & 0.8238 & 0.0160 & 0.8127 & 0.8422 \\
$5\times10^{-4}$ & 0.8158 & 0.0144 & 0.8042 & 0.8319 \\
$7\times10^{-4}$ & 0.8117 & 0.0163 & 0.7994 & 0.8301 \\
$1\times10^{-3}$ & 0.8160 & 0.0170 & 0.7998 & 0.8337 \\
\bottomrule
\end{tabular}
\end{table*}

The defensible conclusion from this grid is limited to the historical execution path: optimizer stabilization widened its usable learning-rate range. It does not establish stability, competitiveness, or predictive superiority of the repaired canonical TransnnMIL. The completed frozen repaired-architecture matched result below governs comparative claims, rather than the historical stability grid.

\section{Retrospective scanner-objective mechanism analysis}
On 14 previously exposed SCORPION TRAIN-AUDIT slides, an independent replay of five original fold-0 seeds found lower standardized feature-space source, target, and observed scanner-contrast reconstruction errors for the \texttt{no\_adversary} variant than for the original full objective (source MSE 0.063253 versus 0.066162, target 0.065253 versus 0.068403, contrast 0.068017 versus 0.072047). Its linear B256 scanner probe also showed lower accuracy. However, changes in pooled scanner centroids did not agree in direction with slide-conditioned centroid variation. These are retrospective mechanistic measurements on previously exposed data: they neither negate the prospective SCORPION registered comparison nor establish an independently validated biological/acquisition split.

\section{Whole-slide and institutional result pointers}
The full PANDA feature manifest contains 10,611 readable Phikon slide bags. TransnnMIL optimization-stability results are recorded under \codepath{docs/results/panda-transnnmil-stability.md}. The full-cache dominance-aware detector experiment is recorded under \codepath{docs/results/dominance-aware-switch-full-panda.md}, and the fixed label-noise-to-ordinal-shift transfer is recorded under \codepath{docs/research/dominance-detector-transfer-results.md}.

CAMELYON17 source-weighting and center-subspace experiments are recorded under \codepath{results/camelyon17/}, \codepath{results/camelyon17_supervised_resnet18/}, and the corresponding research notes.

\section{Reproducibility}
The repository preserves code, tests, result summaries, release manifests, and source hashes for the major research lines. Detailed correction history and superseded analyses remain available in the repository but are supporting research records rather than the narrative focus of the flagship manuscript.

\end{document}
